\documentclass[12pt]{article}
\usepackage[utf8]{inputenc}
\usepackage[margin=0.75in]{geometry}
\usepackage{amsmath}
\usepackage{amssymb}
\usepackage{amsthm}
\usepackage{graphicx} % Required for inserting images
\usepackage{array}
\usepackage{tikz-cd}


\title{MAT1100 Homework 4}
\author{Aaron Avram}
\date{November 17 2026}

\begin{document}
\newtheorem{definition}{Definition}
\newtheorem{theorem}{Theorem}
\newtheorem{claim}{Claim}
\newtheorem{corollary}{Corollary}

\maketitle

\begin{align*}
    \textbf{Question 4}
\end{align*}

\begin{enumerate}
    \item[(i)]
    \begin{proof}
        If $g = 0$ there is nothing to prove, so suppose $g \neq 0$. If $f \in R[x]$ is primitive and $g \in R[x]$.
        Suppose $fh = g$ for some $h \in K[x]$. Now write
        \[ h(x) = \sum_{i = 0}^n a_ix^i \]
        where $a_i = p_i/q_i$ for $p_i, q_i \in R$ with $\gcd(p_i, q_i) = 1$.
        Then let $m = \operatorname{lcm}(q_0, \ldots, q_n)$. Clearly then $mh(x)$ is a polynomial with coefficients in $R$,
        suppose that $a$ is the gcd of its coefficients. Then $\frac{m}{a}h(x)$ is a primitive polynomial
        and
        \[ g = \frac{a}{m}f(\frac{m}{a}h) \]
        now $P := f(\frac{m}{a}h)$ is a primitive polynomial by Gauss's lemma so this implies
        and we have
        \[ mg = aP \]
        and since $P$ is primitive the the gcd of the coefficients of $aP$ is $a$
        but $aP = mg$ so the gcd of its coefficients is divisible by $m$ therefore $\frac{a}{m} \in R$
        and therefore
        \[ \frac{a}{m}(\frac{m}{a}h) = h \in R[x] \]
        as required.
    \end{proof}

    \item[(ii)]
    \begin{proof}
        Suppose $m \in R$
        such that $mh(x) \in R[x]$ and moreover suppose $a$ is the gcd of the coefficients
        of $mh(x)$. Then
        \[ f = (\frac{a}{m}g)(\frac{m}{a}h) \]
        and by part (i) this implies that $\frac{a}{m}g \in R[x]$.
        Now $g, h$ are monic so for $\frac{m}{a}h$ to be in $R[x]$ we must have $a \mid m$
        and similarly for $\frac{a}{m}g$ to be in $R[x]$ we must have $m \mid a$, therefore
        there exists a unit $v$ such that $a = mv$ so
        \[ \frac{a}{m}g = vg \in R[x] \]
        and then multiplying by $v^{-1}$ we find that $g \in R[x]$ and similarly with $h$.
    \end{proof}

    \item[(iii)]
    \begin{proof}
        Suppose $f(x) = \alpha^n + a_{n-1}\alpha^{n-1} + \ldots + a_1\alpha + a_0 = 0$.
        Then $(x - \alpha)$ is a monic factor of $f(x)$. Suppose $g(x) \in K[x]$
        such that $f(x) = (x - \alpha)g(x)$ with $g(x)$ clearly monic as $f(x)$ is monic.
        By (ii) this implies that $(x - \alpha), g(x) \in R[x]$ and therefore $\alpha \in R$.

    \end{proof}
\end{enumerate}
\end{document}